\subsection{Implementación de recomendaciones y fichas de práctica}
\label{sec:impl-recomendaciones}

Al guardar una calificación, \texttt{dao\_tutor.py} actualiza el intento y los análisis confirmados. A continuación, \texttt{dao\_retroalimentacion.py} consulta los promedios del alumno y llama al servicio \texttt{retroalimentacion.py}. El servicio aplica la regla RN-26 y devuelve, si corresponde, el contenido de la práctica. El DAO guarda ese contenido antes de confirmar la transacción. Si falla el guardado de la recomendación, se revierte también la calificación.

La ruta de consulta del alumno recupera las recomendaciones de sus intentos calificados. En el panel del tutor, la respuesta de calificación y el detalle del intento incluyen la recomendación disponible. Ambas vistas presentan la información con el componente \texttt{RetroalimentacionEjercicio.jsx}. Para los registros en formato anterior, el servicio adapta el contenido al formato de las fichas actuales al consultarlo y conserva el motivo y el promedio originales.

\subsubsection{Generación y descarga de las fichas PDF}

Al solicitar la descarga, el servidor verifica que el intento pertenezca al alumno autenticado o a un alumno a cargo del tutor autenticado. El servicio \texttt{ficha\_retroalimentacion.py} construye el PDF en memoria con ReportLab a partir de la recomendación recuperada. El documento incluye los pasos y el material de práctica. Las planas llevan renglones para copiar las palabras del modelo y escribir una oración. Las fichas de motricidad incluyen el laberinto o las guías de papel que corresponden a la actividad.

El servidor entrega un archivo \texttt{practica-<id>.pdf} cuando se solicita. La base de datos conserva la descripción y el modelo de la actividad para reproducir la ficha en descargas posteriores. El alumno puede seguir las instrucciones en pantalla o imprimir el archivo para practicar.

\subsubsection{Registro y validación de la práctica}

Cuando el alumno envía la foto, el servidor valida el archivo y comprueba que la recomendación pertenezca a su intento. Guarda la imagen junto con la fecha de entrega y cambia el estado a \texttt{completada}. Si el tutor vuelve a calificar el intento, se conserva la evidencia ya recibida. Los reintentos de la misma solicitud conservan el registro original.

El tutor solicita la imagen desde el detalle del intento. El servidor comprueba que el alumno esté a su cargo antes de mostrar la evidencia y antes de guardar la validación. Al validar, se registran el tutor, la fecha y el estado \texttt{validada}. La foto se guarda como evidencia de la práctica, sin crear un intento adicional ni ejecutar Cloud Vision.
